\documentclass[11pt,letterpaper]{article}

\usepackage[margin=1in]{geometry}
\usepackage{amsmath,amssymb,bm}
\usepackage{siunitx}
\usepackage{booktabs}
\usepackage{graphicx}
\usepackage{xcolor}
\usepackage{hyperref}
\usepackage{enumitem}
\usepackage{physics}
\usepackage{mathtools}
\usepackage{microtype}

\hypersetup{
    colorlinks=true,
    linkcolor=blue,
    citecolor=blue,
    urlcolor=blue
}

\newcommand{\gtt}{g^{(2)}}
\newcommand{\Gnorm}{G^{(2)}}
\newcommand{\avg}[1]{\left\langle #1 \right\rangle}
\renewcommand{\dd}{\mathrm{d}}
\newcommand{\ii}{\mathrm{i}}

\title{Initial Scaling Analysis of Dynamic Speckle from a Rotating Diffuser\\
and Possible Applications in Information Processing and Random Thermal Excitation}
\author{}
\date{\today}

\begin{document}

\maketitle

\begin{abstract}
This document defines an initial analysis framework for intensity
autocorrelation data obtained from coherent light scattered by a rotating
multiple-scattering diffuser. The measured data consist of second-order
intensity correlation curves at several scattering angles and rotation
speeds. The first analysis goal is to test whether the temporal correlation
curves collapse when time delay is replaced by diffuser angular displacement,
$\Delta\theta=\omega\tau$. A stronger analysis incorporates the tangential
displacement $R\omega\tau$ and a possible scattering-vector-dependent
decorrelation length $\ell_c(q)$. The second part outlines possible
applications in information transmission, random-feature optical computing,
and spatially random thermal excitation for transport measurements. The
document is intended both as a physics note and as a specification for
analysis code.
\end{abstract}

\tableofcontents

\section{Experimental variables and notation}

A coherent optical beam illuminates a rotating diffuser at radial distance
$R$ from the axis of rotation. The diffuser rotates at angular speed
$\omega$. A detector measures scattered intensity at scattering angle
$\phi$, defined as the angle between the incident and detected wavevectors.

The magnitude of the scattering vector is

\begin{equation}
    \boxed{
    q = \frac{4\pi}{\lambda}\sin\left(\frac{\phi}{2}\right)
    }
    \label{eq:q_definition}
\end{equation}

where $\lambda$ is the optical wavelength in the surrounding medium. If the
beam propagates in a material with refractive index $n$, replace $\lambda$
by the wavelength in the medium, $\lambda_0/n$, or write

\begin{equation}
    q = \frac{4\pi n}{\lambda_0}
    \sin\left(\frac{\phi}{2}\right).
\end{equation}

The primary measured quantity is the intensity time series

\begin{equation}
    I(t,q;\omega,R),
\end{equation}

or, when the diffuser angle is measured,

\begin{equation}
    I(\theta,q),
    \qquad
    \theta(t)=\theta_0+\omega t
\end{equation}

for uniform rotation.

The present data consist primarily of normalized second-order intensity
autorrelation functions rather than complete intensity traces.

\section{Definition of the intensity autocorrelation}

For a detector at scattering vector $q$, define

\begin{equation}
    \gtt(\tau,q;\omega,R)
    =
    \frac{
    \avg{
    I(t,q;\omega,R)
    I(t+\tau,q;\omega,R)
    }_t
    }{
    \avg{I(t,q;\omega,R)}_t^2
    }.
    \label{eq:g2_definition}
\end{equation}

For a stationary intensity process,

\begin{equation}
    \gtt(\tau,q;\omega,R)
    \rightarrow 1
\end{equation}

at delays much longer than the correlation time.

It is often useful to remove differences in measured speckle contrast by
defining

\begin{equation}
    \boxed{
    \Gnorm(\tau,q;\omega,R)
    =
    \frac{
    \gtt(\tau,q;\omega,R)-1
    }{
    \gtt(0,q;\omega,R)-1
    }.
    }
    \label{eq:G_normalized}
\end{equation}

Ideally,

\begin{equation}
    \Gnorm(0,q;\omega,R)=1,
    \qquad
    \Gnorm(\tau\rightarrow\infty,q;\omega,R)=0.
\end{equation}

The normalization in Eq.~\eqref{eq:G_normalized} removes the contrast
factor from the comparison of decay shapes. The raw value
$\gtt(0)-1$ should nevertheless be retained and analyzed separately because
it contains information about detector area, polarization averaging,
spatial coherence, and the number of unresolved speckle modes.

\section{Candidate functional form}

A useful empirical model is the stretched or compressed exponential

\begin{equation}
    \gtt(\tau,q;\omega,R)-1
    =
    \beta(q,\omega,R)
    \exp\left[
    -2
    \left|
    \frac{\tau}{\tau_c(q,\omega,R)}
    \right|^\alpha
    \right],
    \label{eq:stretched_g2}
\end{equation}

where

\begin{itemize}[leftmargin=2em]
    \item $\beta=\gtt(0)-1$ is the measured correlation contrast,
    \item $\tau_c$ is a characteristic correlation time,
    \item $\alpha$ controls the decay shape.
\end{itemize}

The factor of 2 in Eq.~\eqref{eq:stretched_g2} is conventional when the
intensity correlation is related to the square of a field correlation.
The analysis code should allow this factor to be removed or absorbed into
the definition of $\tau_c$.

The exponent should initially remain a fitted parameter. Common limiting
cases include

\begin{align}
    \alpha &= 1
    &&\text{exponential decay},\\
    \alpha &= 2
    &&\text{Gaussian decay}.
\end{align}

The main scaling tests should not rely exclusively on fitted values of
$\tau_c$. Full-curve collapse tests use more information and can reveal
changes in shape, shoulders, oscillations, or multiple timescales.

\section{First prediction: collapse versus \texorpdfstring{$\omega\tau$}{omega tau}}

\subsection{Angular displacement as the governing variable}

During a delay $\tau$, a uniformly rotating diffuser advances through angle

\begin{equation}
    \Delta\theta = \omega\tau.
    \label{eq:angular_displacement}
\end{equation}

If rotation speed changes only the rate at which a fixed sequence of diffuser
configurations is traversed, then

\begin{equation}
    I(t,q;\omega,R)
    =
    I\left(\theta_0+\omega t,q;R\right).
\end{equation}

The normalized correlation should then have the form

\begin{equation}
    \boxed{
    \Gnorm(\tau,q;\omega,R)
    =
    F_{q,R}(\omega\tau).
    }
    \label{eq:first_collapse}
\end{equation}

At fixed $q$ and fixed $R$, curves measured at different $\omega$ should
collapse when plotted against

\begin{equation}
    u_\theta=\omega\tau.
\end{equation}

The correlation time must consequently scale as

\begin{equation}
    \boxed{
    \tau_c(q,\omega,R)
    =
    \frac{\Delta\theta_c(q,R)}{\omega},
    }
    \label{eq:tau_omega}
\end{equation}

where $\Delta\theta_c(q,R)$ is the characteristic angular displacement
required to decorrelate the detected intensity.

Thus, at fixed $q$ and $R$,

\begin{equation}
    \tau_c \propto \omega^{-1}.
\end{equation}

A plot of $\tau_c$ versus $1/\omega$ should be approximately linear and,
within uncertainty, should pass through the origin unless there is an
additional speed-independent decorrelation process.

\subsection{What successful collapse means}

Successful collapse versus $\omega\tau$ supports the following interpretation:

\begin{quote}
The normalized speckle dynamics are governed primarily by diffuser angular
displacement. Rotation speed changes the temporal playback rate but does not
substantially change the underlying normalized stochastic process.
\end{quote}

This is stronger than finding only $\tau_c\propto 1/\omega$, because collapse
tests the entire shape of every correlation curve.

It does not prove that intensity traces at different $q$ values are identical
or mutually predictable. It establishes a scaling law for their normalized
second-order statistics.

\subsection{Failure of the first collapse}

Systematic failure of collapse may indicate

\begin{itemize}[leftmargin=2em]
    \item speed-dependent vibration or shaft wobble,
    \item nonuniform angular velocity,
    \item mechanical resonances,
    \item detector bandwidth limitations,
    \item time-bin or correlator artifacts,
    \item changing beam position on the diffuser,
    \item multiple decorrelation processes,
    \item diffuser heating or deformation,
    \item insufficient statistical averaging.
\end{itemize}

Failure is therefore diagnostically useful.

\section{Second prediction: collapse versus tangential displacement}

The material under the illuminated spot moves tangentially through distance

\begin{equation}
    \Delta x = R\Delta\theta=R\omega\tau.
    \label{eq:tangential_displacement}
\end{equation}

If decorrelation is governed by physical translation of the scattering
structure through the illuminated region, the stronger scaling hypothesis is

\begin{equation}
    \boxed{
    \Gnorm(\tau,q;\omega,R)
    =
    F_q(R\omega\tau).
    }
    \label{eq:spatial_collapse}
\end{equation}

Introduce a characteristic physical decorrelation length $\ell_c(q)$:

\begin{equation}
    \boxed{
    \Gnorm(\tau,q;\omega,R)
    =
    F\left(
    \frac{R\omega\tau}{\ell_c(q)}
    \right).
    }
    \label{eq:general_collapse}
\end{equation}

The corresponding correlation time is

\begin{equation}
    \boxed{
    \tau_c(q,\omega,R)
    =
    \frac{\ell_c(q)}{R\omega}.
    }
    \label{eq:tau_general}
\end{equation}

Therefore,

\begin{equation}
    R\omega\tau_c=\ell_c(q).
    \label{eq:decorrelation_length}
\end{equation}

Equation~\eqref{eq:decorrelation_length} provides an operational way to
extract a physical displacement scale from each measured correlation curve.

At fixed $q$, the following predictions can be tested:

\begin{align}
    \tau_c &\propto \omega^{-1},\\
    \tau_c &\propto R^{-1},\\
    R\omega\tau_c &\approx \text{constant}.
\end{align}

\section{Possible dependence on scattering vector}

The most conservative analysis does not assume a particular $q$ dependence.
Instead, determine

\begin{equation}
    \ell_c(q)=R\omega\tau_c(q,\omega,R)
    \label{eq:ellq_empirical}
\end{equation}

from the data and then examine the functional dependence of $\ell_c$ on $q$.

The general collapse variable is

\begin{equation}
    \boxed{
    u=
    \frac{R\omega\tau}{\ell_c(q)}.
    }
    \label{eq:master_variable}
\end{equation}

Candidate models include the following.

\subsection{No measurable \texorpdfstring{$q$}{q} dependence}

If

\begin{equation}
    \ell_c(q)\approx \ell_0,
\end{equation}

then

\begin{equation}
    \Gnorm
    =
    F\left(
    \frac{R\omega\tau}{\ell_0}
    \right).
\end{equation}

All scattering channels have the same physical decorrelation displacement,
although their actual intensity traces may still be different.

\subsection{Inverse-\texorpdfstring{$q$}{q} displacement scale}

A phase-sensitivity argument may suggest

\begin{equation}
    \ell_c(q)\propto \frac{1}{q}.
\end{equation}

Then the dimensionless scaling variable is

\begin{equation}
    u_q=qR\omega\tau,
    \label{eq:q_scaling}
\end{equation}

and

\begin{equation}
    \tau_c\propto \frac{1}{qR\omega}.
\end{equation}

This model should be treated as a hypothesis, not assumed in advance,
especially for a multiple-scattering diffuser whose dynamics may not reduce
to a single-scattering phase factor.

\subsection{General power law}

A more flexible model is

\begin{equation}
    \ell_c(q)
    =
    \ell_0
    \left(\frac{q}{q_0}\right)^{-\nu},
    \label{eq:q_power_law}
\end{equation}

which gives

\begin{equation}
    \tau_c(q,\omega,R)
    =
    \frac{\ell_0}{R\omega}
    \left(\frac{q}{q_0}\right)^{-\nu}.
\end{equation}

The master variable becomes

\begin{equation}
    u
    =
    \frac{R\omega\tau}{\ell_0}
    \left(\frac{q}{q_0}\right)^\nu.
\end{equation}

The reference value $q_0$ keeps the expression dimensionally consistent.
The exponent $\nu$ should be estimated from the data.

\section{Revolution-synchronized recurrence}

The first-collapse prediction concerns normalized statistics. A distinct
question is whether the detailed intensity pattern repeats after one
revolution.

The revolution period is

\begin{equation}
    T_{\mathrm{rev}}=\frac{2\pi}{\omega}.
\end{equation}

Perfect recurrence would imply

\begin{equation}
    I(t+T_{\mathrm{rev}},q)
    =
    I(t,q),
\end{equation}

or equivalently,

\begin{equation}
    I(\theta+2\pi,q)=I(\theta,q).
\end{equation}

Experimentally one expects approximate recurrence,

\begin{equation}
    I(\theta+2\pi,q)\approx I(\theta,q),
\end{equation}

limited by photon statistics, laser noise, diffuser runout, mechanical drift,
and imperfect angle registration.

A useful revolution-to-revolution recurrence coefficient is

\begin{equation}
    \rho_{\mathrm{rev}}(q)
    =
    \frac{
    \avg{
    \delta I_r(\theta,q)
    \delta I_{r+1}(\theta,q)
    }_\theta
    }{
    \sqrt{
    \avg{\delta I_r^2}_\theta
    \avg{\delta I_{r+1}^2}_\theta
    }
    },
    \label{eq:rev_correlation}
\end{equation}

where $r$ labels revolutions.

Collapse and recurrence answer different questions:

\begin{itemize}[leftmargin=2em]
    \item Collapse tests whether the normalized statistics scale with
    angular or tangential displacement.
    \item Recurrence tests whether the detailed diffuser-indexed intensity
    pattern can be replayed and calibrated.
\end{itemize}

A system may show good collapse but poor recurrence. It would then provide
predictable stochastic statistics but not a stable, deterministic sequence
of optical features.

\section{Recommended initial analysis workflow}

The following sequence is intended as a direct guide for analysis code.

\subsection{Input data}

For each experimental run, store

\begin{itemize}[leftmargin=2em]
    \item delay values $\tau_i$,
    \item measured $\gtt(\tau_i)$,
    \item uncertainty or number of averages, when available,
    \item scattering angle $\phi$,
    \item wavelength $\lambda$,
    \item calculated $q$ from Eq.~\eqref{eq:q_definition},
    \item angular speed $\omega$,
    \item beam radius $R$ on the diffuser,
    \item detector aperture and optical configuration,
    \item acquisition duration,
    \item count rate or mean intensity,
    \item rotation direction.
\end{itemize}

Use SI units internally:

\begin{equation}
    q\,[\si{\per\meter}],
    \quad
    R\,[\si{\meter}],
    \quad
    \omega\,[\si{\radian\per\second}],
    \quad
    \tau\,[\si{\second}].
\end{equation}

\subsection{Preprocessing}

For each curve:

\begin{enumerate}[leftmargin=2em]
    \item Estimate the long-delay baseline.
    \item Correct the baseline to unity if justified.
    \item Calculate $\beta=\gtt(0)-1$.
    \item Calculate $\Gnorm$ using Eq.~\eqref{eq:G_normalized}.
    \item Preserve both raw and normalized curves.
    \item Flag curves that do not reach a stable long-delay baseline.
\end{enumerate}

Do not force the final measured point to equal one. Estimate the baseline
from a delay region that appears statistically stationary.

\subsection{Individual-curve fits}

Fit each curve to Eq.~\eqref{eq:stretched_g2}, initially allowing

\begin{equation}
    \beta,\quad \tau_c,\quad \alpha,\quad \text{and baseline}
\end{equation}

to vary.

Also compare restricted models with $\alpha=1$ and $\alpha=2$.

Report

\begin{itemize}[leftmargin=2em]
    \item fitted parameters,
    \item parameter uncertainties,
    \item residuals,
    \item reduced chi-squared or another goodness-of-fit measure,
    \item information criteria for competing models,
    \item the delay range included in the fit.
\end{itemize}

\subsection{First-collapse plot}

For each fixed $q$ and $R$, plot

\begin{equation}
    \Gnorm
    \quad \text{versus} \quad
    \omega\tau.
\end{equation}

The plotting code should show all rotation speeds on the same axes.

Quantify collapse error by interpolating the curves onto a shared
dimensionless grid and calculating, for example,

\begin{equation}
    \chi^2_{\mathrm{collapse}}
    =
    \frac{1}{N_{\mathrm{tot}}}
    \sum_{j,i}
    \frac{
    \left[
    \Gnorm_j(u_i)-\overline{\Gnorm}(u_i)
    \right]^2
    }{
    \sigma_{j,i}^2
    },
    \label{eq:collapse_error}
\end{equation}

where $j$ labels runs and $\overline{\Gnorm}$ is the estimated master curve.

If reliable pointwise uncertainties are unavailable, use an unweighted RMS
collapse error and estimate uncertainty by bootstrap resampling.

\subsection{Correlation-time scaling}

At fixed $q$ and $R$, fit

\begin{equation}
    \tau_c=\frac{A}{\omega}+B.
\end{equation}

The ideal angular-displacement model predicts

\begin{equation}
    B=0,
    \qquad
    A=\Delta\theta_c.
\end{equation}

A statistically significant nonzero $B$ may indicate an additional
timescale or a systematic analysis offset.

\subsection{Tangential-displacement collapse}

When multiple $R$ values are available, plot

\begin{equation}
    \Gnorm
    \quad \text{versus} \quad
    R\omega\tau.
\end{equation}

Calculate

\begin{equation}
    \ell_c(q)=R\omega\tau_c
\end{equation}

for each run. At fixed $q$, test whether the extracted values agree within
uncertainty.

\subsection{\texorpdfstring{$q$}{q}-dependent collapse}

Plot

\begin{equation}
    \ell_c(q)=R\omega\tau_c
\end{equation}

against $q$. Compare at least

\begin{align}
    \ell_c(q)&=\ell_0,\\
    \ell_c(q)&=\frac{A}{q},\\
    \ell_c(q)&=\ell_0\left(\frac{q}{q_0}\right)^{-\nu}.
\end{align}

Only after estimating $\ell_c(q)$ should all data be plotted against

\begin{equation}
    u=\frac{R\omega\tau}{\ell_c(q)}.
\end{equation}

The strongest collapse claim would be

\begin{equation}
    \boxed{
    \Gnorm(\tau,q;\omega,R)
    =
    F\left(
    \frac{R\omega\tau}{\ell_c(q)}
    \right)
    }
\end{equation}

over the measured ranges of $q$, $\omega$, and $R$.

\subsection{Useful plots}

The initial analysis package should generate:

\begin{enumerate}[leftmargin=2em]
    \item Raw $\gtt(\tau)$ curves grouped by $q$.
    \item Normalized $\Gnorm(\tau)$ curves.
    \item Fits and fit residuals.
    \item $\tau_c$ versus $1/\omega$ at fixed $q$ and $R$.
    \item $\Gnorm$ versus $\omega\tau$.
    \item $\Gnorm$ versus $R\omega\tau$, when $R$ varies.
    \item $\ell_c(q)=R\omega\tau_c$ versus $q$.
    \item Global collapse versus $R\omega\tau/\ell_c(q)$.
    \item $\alpha$ versus $q$, $\omega$, and $R$.
    \item $\beta$ versus $q$, $\omega$, and $R$.
    \item Collapse residuals versus scaled delay.
\end{enumerate}

\section{Autocorrelation, spectra, and information}

Define the centered intensity fluctuation

\begin{equation}
    \delta I(t,q)=I(t,q)-\avg{I(t,q)}.
\end{equation}

Its autocovariance is

\begin{equation}
    C_{qq}(\tau)
    =
    \avg{
    \delta I(t,q)
    \delta I(t+\tau,q)
    }.
\end{equation}

It is related to $\gtt$ by

\begin{equation}
    C_{qq}(\tau)
    =
    \avg{I(q)}^2
    \left[
    \gtt(\tau,q)-1
    \right].
\end{equation}

By the Wiener--Khinchin theorem, the Fourier transform of the
autocovariance is the intensity-fluctuation power spectral density:

\begin{equation}
    S_{qq}(f)
    =
    \int_{-\infty}^{\infty}
    C_{qq}(\tau)
    e^{-\ii 2\pi f\tau}
    \dd\tau.
    \label{eq:wiener_khinchin}
\end{equation}

Thus $\gtt(\tau,q)$ determines the spectral power at that detector, apart
from normalization and the DC contribution.

However, it does not determine the complex Fourier phase of the original
intensity trace. Many different intensity realizations can have the same
autocorrelation and power spectrum.

Consequently, universal autocorrelation scaling can coexist with a large
information-bearing space:

\begin{quote}
The family of signals may have a simple low-dimensional statistical
description while the individual intensity realizations remain complex,
distinct, and potentially high dimensional.
\end{quote}

\section{Multiple detectors and cross-correlation}

For detectors at $q_1$ and $q_2$, define

\begin{align}
    I_1(t)&=I(t,q_1),\\
    I_2(t)&=I(t,q_2).
\end{align}

The cross-covariance is

\begin{equation}
    C_{12}(\tau)
    =
    \avg{
    \delta I_1(t)
    \delta I_2(t+\tau)
    }.
    \label{eq:cross_covariance}
\end{equation}

The normalized cross-correlation is

\begin{equation}
    \rho_{12}(\tau)
    =
    \frac{
    C_{12}(\tau)
    }{
    \sqrt{
    C_{11}(0)C_{22}(0)
    }
    }.
    \label{eq:normalized_cross}
\end{equation}

The Fourier transform gives the cross-spectrum

\begin{equation}
    S_{12}(f)
    =
    \int_{-\infty}^{\infty}
    C_{12}(\tau)e^{-\ii2\pi f\tau}\dd\tau.
\end{equation}

Unlike an autospectrum, $S_{12}$ is generally complex. Its phase describes
the relative spectral phase between the two detector signals.

A useful normalized quantity is the magnitude-squared coherence,

\begin{equation}
    \gamma_{12}^2(f)
    =
    \frac{
    |S_{12}(f)|^2
    }{
    S_{11}(f)S_{22}(f)
    },
    \qquad
    0\leq\gamma_{12}^2\leq 1.
    \label{eq:coherence}
\end{equation}

Low coherence indicates that the detectors carry largely distinct
second-order information at that frequency. High coherence indicates strong
redundancy or a common physical component.

Two autocorrelation functions may collapse to the same master form after
rescaling even while the detector signals themselves are not identical.
Collapse concerns the form of the statistics. Cross-correlation concerns
the shared information between channels.

\section{Lag direction and causal interpretation}

The cross-correlation in Eq.~\eqref{eq:cross_covariance} is directional in
its lag convention:

\begin{equation}
    C_{12}(\tau)
    =
    \avg{
    \delta I_1(t)
    \delta I_2(t+\tau)
    }.
\end{equation}

A peak at positive $\tau$ means that variations at detector 1 tend to
precede similar variations at detector 2 by $\tau$. A peak at negative
$\tau$ means the opposite ordering.

For example, if

\begin{equation}
    I_2(t)\approx aI_1(t-\tau_d)+\eta(t),
\end{equation}

then the cross-correlation will generally peak near $\tau=\tau_d$, subject
to the sign convention.

This may reveal

\begin{itemize}[leftmargin=2em]
    \item propagation delays,
    \item speckle translation across detector positions,
    \item ordered traversal of scattering channels,
    \item thermal transport delays,
    \item detector or electronics delays,
    \item common-input responses with different transfer functions.
\end{itemize}

A nonzero-lag peak does not by itself prove physical causality. Common
drivers, filtering, periodicity, or geometric motion can produce similar
signatures. A causal claim requires a controlled perturbation, known timing,
or an explicit dynamical model.

With a driven system, it is useful to correlate each detector with the
known drive $u(t)$:

\begin{equation}
    C_{uI_m}(\tau)
    =
    \avg{
    \delta u(t)
    \delta I_m(t+\tau)
    }.
    \label{eq:input_output_cross}
\end{equation}

The known drive supplies a causal time origin. The sign and shape of
$C_{uI_m}$ can then be interpreted as an input--output response more
reliably than the detector--detector correlation alone.

\section{Part II: Motivation and possible applications}

\section{Random optical feature generation}

Let an optical input be represented by a vector $\bm{x}$. It may encode

\begin{itemize}[leftmargin=2em]
    \item amplitudes in several spatial modes,
    \item an image or spatial mask,
    \item phases imposed by an SLM,
    \item several wavelength channels,
    \item polarization components,
    \item several temporal input channels.
\end{itemize}

For a diffuser angle $\theta$ and output channel $q_m$, a schematic field
model is

\begin{equation}
    E(\theta,q_m;\bm{x})
    =
    \sum_j T_{mj}(\theta)x_j,
\end{equation}

where $T_{mj}(\theta)$ is an angle-dependent optical transfer coefficient.
The measured intensity is

\begin{equation}
    I(\theta,q_m;\bm{x})
    =
    \left|
    \sum_j T_{mj}(\theta)x_j
    \right|^2.
    \label{eq:intensity_feature}
\end{equation}

Linear optical propagation followed by square-law detection produces a
nonlinear feature map of the input.

A finite feature vector may be assembled from several detector channels and
several diffuser angles:

\begin{equation}
    \bm{f}(\bm{x})
    =
    \left\{
    I(\theta_n,q_m;\bm{x})
    \right\}_{n,m}.
    \label{eq:feature_vector}
\end{equation}

A trained linear readout can then form

\begin{equation}
    \widehat{\bm{y}}
    =
    W\bm{f}(\bm{x})+\bm{b}.
\end{equation}

The diffuser performs complicated optical mixing. Only the electronic
readout needs to be trained.

\subsection{Why collapse matters for computing}

If

\begin{equation}
    \Gnorm(\tau,q;\omega,R)
    =
    F\left(
    \frac{R\omega\tau}{\ell_c(q)}
    \right),
\end{equation}

then the normalized feature dynamics are governed by a predictable scaling
law. The feature-refresh time is approximately

\begin{equation}
    \tau_c(q)
    =
    \frac{\ell_c(q)}{R\omega}.
\end{equation}

The approximate feature-refresh rate is therefore

\begin{equation}
    \mathcal{R}_f(q)
    \sim
    \frac{1}{\tau_c(q)}
    =
    \frac{R\omega}{\ell_c(q)}.
\end{equation}

Changing $\omega$ changes the rate at which the diffuser configurations are
traversed. If recurrence is strong, the device resembles a cyclic bank of
reproducible random transformations indexed by diffuser angle.

Collapse does not prove high computational dimension. That requires
measurements across multiple angles, detectors, and inputs, followed by
covariance, singular-value, mutual-information, or task-performance
analysis.

\section{Possible application to information transmission}

A controlled diffuser can produce a reproducible, complicated carrier
indexed by $\theta$ and $q$. Information could be encoded in the incident
field $\bm{x}$ and recovered from one or more scattered intensity streams.

Possible architectures include:

\begin{enumerate}[leftmargin=2em]
    \item \textbf{Calibrated random projection.}
    The receiver learns the map between input symbols and output intensity
    features.

    \item \textbf{Spread-spectrum-like encoding.}
    An input symbol is distributed over many angular positions, detector
    channels, or spectral components.

    \item \textbf{Challenge--response transmission.}
    The diffuser configuration acts as a physical key or challenge variable.

    \item \textbf{Low-detector-count reception.}
    One or a few photon-counting detectors acquire sequential angular
    features rather than a full camera image.

    \item \textbf{Joint spatial and temporal multiplexing.}
    Several $q$ channels are detected in parallel while rotation supplies
    sequential features.
\end{enumerate}

The information rate cannot be inferred from $\tau_c$ alone. It also depends
on

\begin{itemize}[leftmargin=2em]
    \item detected photon rate,
    \item shot noise and background,
    \item recurrence fidelity,
    \item cross-channel redundancy,
    \item input alphabet,
    \item calibration drift,
    \item detector bandwidth,
    \item the conditioning of the input-to-feature map.
\end{itemize}

A first communication experiment could compare classification error for
several imposed input patterns as the number of detector channels and
angular samples increases.

\section{Fiber-coupler architectures}

A nominal $2\times2$ fiber coupler has four physical ports. Depending on how
it is used, it can split, combine, or interfere two optical fields.

Several arrangements may be useful.

\subsection{Signal and reference splitting}

One branch can monitor incident laser power while another enters the
diffuser experiment. Correlation with the reference detector can remove or
identify common laser fluctuations.

Let $I_{\mathrm{ref}}(t)$ be the reference signal. Then

\begin{equation}
    C_{\mathrm{ref},m}(\tau)
    =
    \avg{
    \delta I_{\mathrm{ref}}(t)
    \delta I_m(t+\tau)
    }.
\end{equation}

This distinguishes input-source noise from diffuser-generated dynamics.

\subsection{Two scattered channels}

Fibers positioned at $q_1$ and $q_2$ can feed two detectors directly or be
routed through a coupler. Direct detection preserves separate intensity
streams. A coupler can instead form optical sums and differences before
detection.

\subsection{Interferometric combination}

If coherent fields $E_1$ and $E_2$ enter the two input ports of an ideal
balanced coupler, the output intensities contain interference terms:

\begin{align}
    I_+ &\propto |E_1+E_2|^2,\\
    I_- &\propto |E_1-E_2|^2,
\end{align}

up to coupler phase conventions.

Their difference isolates a field-interference quadrature:

\begin{equation}
    I_+-I_-
    \propto
    4\operatorname{Re}\left(E_1E_2^*\right).
\end{equation}

With an imposed phase shift, the complementary quadrature can be measured.
Such a configuration can recover field-sensitive information that is absent
from independent intensity measurements, provided coherence, polarization,
and path stability are maintained.

A simple intensity cross-correlator and an optical interferometer are not
equivalent. The first measures statistical relations between detected
intensities. The second mixes optical fields before square-law detection.

\subsection{Balanced detection}

Two output ports of a balanced coupler can be sent to a balanced detector.
This may suppress common-mode intensity noise and enhance a weak modulated
or interferometric component.

For photon-counting operation, detector dead time, afterpulsing, timing
jitter, and accidental coincidences must be included in the analysis.

\section{Randomly driven thermal fields}

A dynamic speckle field can also serve as a spatially and temporally random
heating pattern. Let the absorbed optical power density be

\begin{equation}
    P(\bm{r},t)
    =
    A(\bm{r})I_{\mathrm{pump}}(\bm{r},t),
\end{equation}

where $A(\bm{r})$ represents optical absorption.

For small perturbations, the temperature response can be written as a
linear convolution,

\begin{equation}
    \Delta T(\bm{r},t)
    =
    \int_{-\infty}^{t}\dd t'
    \int\dd^d r'\,
    \chi_T(\bm{r}-\bm{r}',t-t')
    P(\bm{r}',t'),
    \label{eq:thermal_response}
\end{equation}

where $\chi_T$ is a causal thermal response function.

In Fourier space,

\begin{equation}
    \Delta T(\bm{k},\Omega)
    =
    \chi_T(\bm{k},\Omega)
    P(\bm{k},\Omega).
    \label{eq:thermal_transfer}
\end{equation}

For ordinary Fourier diffusion in a homogeneous material,

\begin{equation}
    \left(
    -\ii\Omega C+\kappa k^2
    \right)
    \Delta T(\bm{k},\Omega)
    =
    P(\bm{k},\Omega),
\end{equation}

so

\begin{equation}
    \chi_T(\bm{k},\Omega)
    =
    \frac{1}{
    \kappa k^2-\ii\Omega C
    },
    \label{eq:diffusive_chi}
\end{equation}

where $C$ is volumetric heat capacity and $\kappa$ is thermal conductivity.

The characteristic diffusive decay time at spatial frequency $k$ is

\begin{equation}
    \tau_{\mathrm{th}}(k)
    =
    \frac{C}{\kappa k^2}
    =
    \frac{1}{Dk^2},
\end{equation}

where $D=\kappa/C$ is thermal diffusivity.

A random heating field excites many $\bm{k}$ and $\Omega$ components
simultaneously. If the incident random field is measured or reproducible,
cross-correlation between the drive and the response can estimate the
thermal transfer function over a range of spatial and temporal frequencies.

\section{Random transient-grating transport measurements}

A conventional transient-grating experiment uses a spatially periodic
excitation with a selected grating wavevector. A random-transient-grating
approach would instead use a known or measured random spatial pattern that
contains a broad distribution of wavevectors.

Let

\begin{equation}
    P(\bm{r},t)
\end{equation}

be the random optical heating field and let

\begin{equation}
    Y_m(t)
\end{equation}

be a probe observable measured at detector $m$. The observable might be

\begin{itemize}[leftmargin=2em]
    \item diffracted probe intensity,
    \item transient reflectivity,
    \item transient transmission,
    \item optical phase,
    \item surface displacement,
    \item x-ray diffuse or Bragg scattering,
    \item acoustic response,
    \item electrical response coupled to temperature.
\end{itemize}

For a linear system,

\begin{equation}
    Y_m(t)
    =
    \int_{-\infty}^{t}
    h_m(t-t')u(t')\dd t'
    +
    \eta_m(t),
    \label{eq:linear_system}
\end{equation}

where $u(t)$ is a measured drive channel and $h_m$ is the causal impulse
response.

The input--output cross-spectrum is

\begin{equation}
    S_{uY_m}(\Omega)
    =
    H_m(\Omega)S_{uu}(\Omega)
\end{equation}

when the additive noise is uncorrelated with the drive. Therefore,

\begin{equation}
    \boxed{
    H_m(\Omega)
    =
    \frac{
    S_{uY_m}(\Omega)
    }{
    S_{uu}(\Omega)
    }
    }
    \label{eq:transfer_estimator}
\end{equation}

wherever the drive spectrum has adequate power and coherence.

This is a direct motivation for recording a reference copy of the random
drive. A four-port fiber-coupler arrangement could send one portion of the
pump-derived signal to a reference detector while another portion drives
the sample.

\subsection{Possible transport regimes}

A broadband random thermal drive could probe:

\begin{itemize}[leftmargin=2em]
    \item ordinary diffusive thermal transport,
    \item multiple thermal diffusion lengths in layered materials,
    \item anisotropic in-plane transport,
    \item interface thermal conductance,
    \item coupled carrier and phonon transport,
    \item thermoelastic acoustic generation,
    \item hydrodynamic or quasiballistic departures from Fourier diffusion,
    \item spatially heterogeneous transport,
    \item nonlinear response as drive amplitude increases.
\end{itemize}

The advantage is not that the random drive creates new thermal physics. The
advantage is that it may excite many spatial and temporal modes in one
measurement and allow them to be separated statistically through known
input correlations.

\subsection{Multiple-detector thermal measurements}

Several response detectors may be placed at different spatial positions,
scattering vectors, or probe geometries. Their signals can be written

\begin{equation}
    \bm{Y}(t)
    =
    \begin{pmatrix}
    Y_1(t)\\
    Y_2(t)\\
    \vdots\\
    Y_M(t)
    \end{pmatrix}.
\end{equation}

The input--output correlation matrix is

\begin{equation}
    \bm{C}_{uY}(\tau)
    =
    \begin{pmatrix}
    C_{uY_1}(\tau)\\
    C_{uY_2}(\tau)\\
    \vdots\\
    C_{uY_M}(\tau)
    \end{pmatrix}.
\end{equation}

The detector--detector matrix is

\begin{equation}
    \bm{C}_{YY}(\tau)
    =
    \left[
    C_{Y_mY_n}(\tau)
    \right]_{m,n}.
\end{equation}

The asymmetry between positive and negative lag may reveal ordered response
or propagation. The known drive is still needed to assign causal direction
reliably.

\section{Relation to fluctuation--dissipation ideas}

The fluctuation--dissipation theorem relates equilibrium fluctuations to
linear response under specific conditions. Schematically, a susceptibility
or dissipative response can be related to an equilibrium correlation
spectrum.

For a classical variable $A$ near thermal equilibrium, one commonly
encounters relations of the form

\begin{equation}
    S_{AA}(\Omega)
    \propto
    \frac{k_{\mathrm B}T}{\Omega}
    \operatorname{Im}\chi_{AA}(\Omega),
    \label{eq:fdt_schematic}
\end{equation}

with the exact prefactors and symmetrization depending on conventions and
on the variables involved.

The proposed rotating-speckle experiment is not automatically an
application of the fluctuation--dissipation theorem. The optical intensity
fluctuations are externally generated, and a laser-heated sample is
generally driven away from equilibrium.

The closer analogy is system identification with stochastic excitation:

\begin{equation}
    \text{known random drive}
    +
    \text{drive--response cross-correlation}
    \longrightarrow
    \text{linear response function}.
\end{equation}

This resembles the logic of fluctuation--dissipation methods because
correlations reveal response functions, but the random drive is imposed
rather than arising from spontaneous thermal equilibrium fluctuations.

The distinction should remain explicit:

\begin{itemize}[leftmargin=2em]
    \item \textbf{Equilibrium FDT:} spontaneous equilibrium fluctuations
    determine linear response.
    \item \textbf{Stochastic system identification:} a measured external
    random drive is used to determine the causal response.
\end{itemize}

If the optical perturbation is weak enough for linear response and the
unperturbed sample remains near equilibrium, the measured response function
can still be compared with equilibrium transport theory. Strong driving,
nonlinear response, or a nonequilibrium steady state requires a more general
framework.

\section{Proposed staged experimental program}

\subsection{Stage 1: scaling of single-detector autocorrelations}

Measure $\gtt(\tau)$ for several

\begin{equation}
    q,\qquad \omega,\qquad R.
\end{equation}

Test, in order,

\begin{align}
    \Gnorm &= F_{q,R}(\omega\tau),\\
    \Gnorm &= F_q(R\omega\tau),\\
    \Gnorm &= F\left(
    \frac{R\omega\tau}{\ell_c(q)}
    \right).
\end{align}

\subsection{Stage 2: recurrence}

Add an angular encoder or index pulse. Test

\begin{equation}
    I(\theta+2\pi,q)\approx I(\theta,q)
\end{equation}

and measure recurrence over many revolutions.

\subsection{Stage 3: two-detector correlations}

Measure

\begin{equation}
    C_{11}(\tau),
    \qquad
    C_{22}(\tau),
    \qquad
    C_{12}(\tau).
\end{equation}

Determine whether the two autocorrelations obey the same scaled master law
and how much independent information remains in the cross-spectrum and
coherence.

\subsection{Stage 4: reference-channel measurement}

Split off a reference signal with a fiber coupler or beamsplitter. Measure

\begin{equation}
    C_{\mathrm{ref},1}(\tau),
    \qquad
    C_{\mathrm{ref},2}(\tau).
\end{equation}

Use these measurements to separate source fluctuations from diffuser and
sample dynamics.

\subsection{Stage 5: controlled input discrimination}

Apply several known optical inputs $\bm{x}$ and test whether features derived
from

\begin{equation}
    I(\theta,q_m;\bm{x})
\end{equation}

improve discrimination or parameter estimation.

Compare

\begin{itemize}[leftmargin=2em]
    \item raw angle-binned intensities,
    \item angular Fourier coefficients,
    \item autocorrelation features,
    \item auto- and cross-correlation features,
    \item reduced features selected by SVD or regularization.
\end{itemize}

\subsection{Stage 6: driven thermal response}

Use the measured or calibrated random optical field as a heating drive.
Measure one or more thermal or structural responses and estimate transfer
functions using input--output cross-correlations.

\section{Primary initial predictions}

The initial data-analysis code should treat the following as explicit
testable predictions.

\begin{enumerate}[leftmargin=2em]
    \item At fixed $q$ and $R$,
    \begin{equation}
        \tau_c\propto\frac{1}{\omega}.
    \end{equation}

    \item At fixed $q$ and $R$, normalized curves collapse as
    \begin{equation}
        \Gnorm(\tau,q;\omega,R)
        =
        F_{q,R}(\omega\tau).
    \end{equation}

    \item If tangential displacement is the governing coordinate,
    \begin{equation}
        \tau_c\propto\frac{1}{R\omega}
    \end{equation}
    and
    \begin{equation}
        \Gnorm
        =
        F_q(R\omega\tau).
    \end{equation}

    \item The empirical decorrelation length is
    \begin{equation}
        \ell_c(q)=R\omega\tau_c.
    \end{equation}

    \item A possible global master law is
    \begin{equation}
        \Gnorm
        =
        F\left(
        \frac{R\omega\tau}{\ell_c(q)}
        \right).
    \end{equation}

    \item A possible, but not assumed, special case is
    \begin{equation}
        \ell_c(q)\propto q^{-1},
    \end{equation}
    leading to
    \begin{equation}
        \Gnorm=F(qR\omega\tau).
    \end{equation}

    \item The fitted decay exponent $\alpha$ should be approximately
    invariant under a true one-parameter collapse. A systematic dependence
    of $\alpha$ on $\omega$, $R$, or $q$ indicates that the curve shape is
    changing and that a simple rescaling is incomplete.

    \item Collapse of autocorrelations does not imply that intensity traces
    at different $q$ are identical. This must be tested with cross-correlation,
    coherence, covariance-rank, or mutual-information measurements.
\end{enumerate}

\section{References}

\begin{thebibliography}{99}

\bibitem{GoodmanSpeckle}
J. W. Goodman,
\textit{Speckle Phenomena in Optics: Theory and Applications},
2nd ed. (SPIE Press, Bellingham, 2020).

\bibitem{Saleh1975}
B. E. A. Saleh, J. Tavolacci, and D. E. Soper,
``Speckle correlation measurement of the velocity of a small rotating
diffuse object,''
\textit{Applied Optics} \textbf{14}, 2344--2349 (1975).
\href{https://doi.org/10.1364/AO.14.002344}
{doi:10.1364/AO.14.002344}.

\bibitem{Takai1980}
N. Takai, T. Iwai, and T. Asakura,
``Real-time velocity measurement for a diffuse object using
zero-crossings of speckle intensity fluctuations,''
\textit{Journal of the Optical Society of America} \textbf{70},
450--455 (1980).
\href{https://doi.org/10.1364/JOSA.70.000450}
{doi:10.1364/JOSA.70.000450}.

\bibitem{Churnside1982}
J. H. Churnside,
``Speckle from a rotating diffuse object,''
\textit{Journal of the Optical Society of America} \textbf{72},
1464--1469 (1982).
\href{https://doi.org/10.1364/JOSA.72.001464}
{doi:10.1364/JOSA.72.001464}.

\bibitem{Yoshimura1986}
T. Yoshimura, K. Kato, and K. Nakagawa,
``Rotational and boiling motion of speckles in a two-lens imaging
system,''
\textit{Journal of the Optical Society of America A} \textbf{3},
1018--1024 (1986).
\href{https://doi.org/10.1364/JOSAA.3.001018}
{doi:10.1364/JOSAA.3.001018}.

\bibitem{Partlo1993}
W. N. Partlo, I. V. Fomenkov, R. M. Ness, and R. I. Sandstrom,
``Diffuser speckle model: application to multiple moving diffuser
laser-beam homogenizers,''
\textit{Applied Optics} \textbf{32}, 3009--3015 (1993).
\href{https://doi.org/10.1364/AO.32.003009}
{doi:10.1364/AO.32.003009}.

\bibitem{Wiener}
N. Wiener,
``Generalized harmonic analysis,''
\textit{Acta Mathematica} \textbf{55}, 117--258 (1930).

\bibitem{Khinchin}
A. Khintchine,
``Korrelationstheorie der station\"aren stochastischen Prozesse,''
\textit{Mathematische Annalen} \textbf{109}, 604--615 (1934).

\bibitem{EichlerTG}
H. J. Eichler, P. G\"unter, and D. W. Pohl,
\textit{Laser-Induced Dynamic Gratings}
(Springer, Berlin, 1986).

\bibitem{JohnsonTG}
J. A. Johnson, A. A. Maznev, J. Cuffe, J. K. Eliason,
A. J. Minnich, T. Kehoe, C. M. Sotomayor Torres, G. Chen,
and K. A. Nelson,
``Direct measurement of room-temperature nondiffusive thermal
transport over micron distances in a silicon membrane,''
\textit{Physical Review Letters} \textbf{110}, 025901 (2013).
\href{https://doi.org/10.1103/PhysRevLett.110.025901}
{doi:10.1103/PhysRevLett.110.025901}.

\bibitem{Song2024}
Q. Song \textit{et al.},
``Probing carrier and phonon transport in semiconductors all at once
through frequency-domain transient grating spectroscopy,''
\textit{Physical Review Applied} \textbf{21}, 034044 (2024).
\href{https://doi.org/10.1103/PhysRevApplied.21.034044}
{doi:10.1103/PhysRevApplied.21.034044}.

\bibitem{Kubo}
R. Kubo,
``The fluctuation-dissipation theorem,''
\textit{Reports on Progress in Physics} \textbf{29}, 255--284 (1966).
\href{https://doi.org/10.1088/0034-4885/29/1/306}
{doi:10.1088/0034-4885/29/1/306}.

\bibitem{BendatPiersol}
J. S. Bendat and A. G. Piersol,
\textit{Random Data: Analysis and Measurement Procedures},
4th ed. (Wiley, Hoboken, 2010).

\bibitem{Ljung}
L. Ljung,
\textit{System Identification: Theory for the User},
2nd ed. (Prentice Hall, Upper Saddle River, 1999).

\end{thebibliography}

\end{document}